\documentclass[12pt]{article}
\usepackage{amsmath}
\usepackage{amsfonts}
\usepackage{amssymb}
\usepackage{geometry}
\geometry{a4paper, margin=1in}
\pagenumbering{gobble}

\title{02YMECH - Homework 1}
\date{Assigned for the week of Sep 22, 2025}
\begin{document}
\maketitle

\section*{Questions}
\begin{enumerate}
    \item Compute following indefinite integrals
          \[
            \int \frac{dx}{a+x} \text{ and } \int \left(3x^2 \, e^{x^3} + \frac{4x}{1 + x^2}\right)\,dx,
          \]
          where \(a\) is a constant.

    \item Compute the definite integral
         \[
           \int_{1}^{4} (2x + 3)\,dx.
         \]
    
    \item Find the derivative of
          \[
            f(x) = \frac{x^3 \ln(x)}{e^{2x}}, \qquad (x>0).
          \]
    
    \item  Find the local minimum and maximum values of the function
           \[
             f(x) = x^3 - 6x^2 + 9x + 1
           \]
           by using the first and second derivative tests. 
    
    \item Find the Taylor series of $\sin x$, $\cos x$, and $e^x$ about $x=0$ up to and including the term in $x^5$.

    \item Compute the partial derivatives of the function
      \[
        f(x, y) = x^2 y + e^{xy}
      \]
      with respect to \(x\) and \(y\) and write down the total differential \(df\).
    
    \item The position of a particle is described by the following vector as a function of time~\(t\):
          \[
        \mathbf{r}(t) = a\cos(t)\,\hat{\imath} + a\sin(t)\,\hat{\jmath} + ct\,\hat{k} 
          \]
          where \(a\) and \(c\) are constants with units of length and velocity, respectively.
          \begin{enumerate}
            \item What kind of trajectory does the particle follow?
            \item Find the velocity, acceleration and speed of the particle as functions of time. 
            \item When you consider the trajectory of the particle again, are the directions of the velocity and acceleration vectors in the way you expect? Explain.
            \item Calculate the projection of the acceleration vector onto the velocity vector as a function of time and give a physical interpretation of this result.
          \end{enumerate}

    \item A flat rectangular surface lies in the $xy$-plane. The surface has sides of length $a = 0.20~\mathrm{m}$ along the $x$-axis and $b = 0.30~\mathrm{m}$ along the $y$-axis. A uniform water flow is described by the flux vector field:

\[
\mathbf{F} = 200\,\hat{\imath} + 300\,\hat{\jmath} + 400\,\hat{k} 
\quad \left[\frac{\mathrm{L}}{\mathrm{m}^2 \cdot \mathrm{s}}\right]
\]

\begin{enumerate}
    \item Compute the area vector $\mathbf{A}$ of the surface (the vector normal to the surface with magnitude equal to the area of the surface).
    \item Compute the volume of water per unit time (in liters per second) passing through the surface using:
    \[
    Q = \mathbf{F} \cdot \mathbf{A}.
    \]
    \item Discuss how the orientation of the surface affects the water flow through it.
\end{enumerate}

    \item A particle moves in a plane such that its position vector at time \(t\) is given by
          \[
            \mathbf{r}(t) = at^2\,\hat{\imath} + bt^3\,\hat{\jmath}.
          \]
          where $a$ and $b$ are constants with appropriate units.
          \begin{enumerate}
            \item Find the velocity and acceleration vectors as functions of time.
            \item Find the angle between the velocity and acceleration vectors at time \(t = 1\).
          \end{enumerate}

    \item Given the vectors
\[
\mathbf{A} = 2\,\hat{\imath} - 3\,\hat{\jmath} + \hat{k}, \quad 
\mathbf{B} = -\hat{\imath} + 4\,\hat{\jmath} + 2\,\hat{k},
\]
compute $\mathbf{A} \times \mathbf{B}$.
    
          
\end{enumerate}
\end{document}
